\section{Testing and Evaluation}\label{ch:impl_testing}

Software that decides about live traffic has to be checked carefully, because a bad
routing or scaling choice is not a cosmetic bug and can take a service down. SmartLoad
therefore treats its tests and benchmarks as engineering work: a fast deterministic
check guards every engine on its own, a slower live-stack check guards every pair of
services on a running cluster, and a small set of end-to-end checks guard each feature
a customer uses. This section lays out that test pyramid, then turns to the custom
benchmark harnesses that answer the project's research questions and the results they
produced.

\subsection{Test strategy: a layered pyramid}\label{imc:pyramid}

The repository organises its checks as a classic test pyramid plus three structural
lints that protect its own conventions. The shape is bottom-heavy on purpose: most
checks are fast, pure-Python unit tests that need no Docker and no network, so the
inner development loop stays in seconds, while expensive live-stack tests are reserved
for contracts that cannot be exercised without a running cluster. \Cref{fig:imc:pyramid}
draws the shape qualitatively, and \cref{tab:imc:layers} reports the per-module
unit-suite counts at the audited commit, the layer where the bulk of the testing
lives.

\begin{figure}[H]
  \centering
  \resizebox{0.95\linewidth}{!}{%
  \begin{tikzpicture}[
    font=\footnotesize,
    layer/.style={draw, align=center, inner sep=3pt, minimum height=9mm},
  ]
    % Pyramid built from trapezoidal layers (widest at the base).
    \fill[slmint!12] (-1.0,0) -- (1.0,0) -- (0.45,1.0) -- (-0.45,1.0) -- cycle;
    \fill[slmint!22] (-2.4,-1.1) -- (2.4,-1.1) -- (1.0,0) -- (-1.0,0) -- cycle;
    \fill[slmint!32] (-4.2,-2.4) -- (4.2,-2.4) -- (2.4,-1.1) -- (-2.4,-1.1) -- cycle;
    \fill[slmint!45] (-6.4,-4.0) -- (6.4,-4.0) -- (4.2,-2.4) -- (-4.2,-2.4) -- cycle;
    \draw (-6.4,-4.0) -- (0,1.0) -- (6.4,-4.0) -- cycle;
    \draw (-4.2,-2.4) -- (4.2,-2.4);
    \draw (-2.4,-1.1) -- (2.4,-1.1);
    \draw (-1.0,0) -- (1.0,0);
    \node at (0,0.55) {\textbf{Conformance}};
    \node[font=\scriptsize] at (0,0.28) {1 suite (nginx + mock)};
    \node at (0,-0.55) {\textbf{End-to-end}};
    \node at (0,-1.75) {\textbf{Integration}};
    \node at (0,-3.1) {\textbf{Unit}};
    % Side annotations on cost / fidelity axis.
    \draw[slflow] (6.9,-4.0) -- node[right, font=\scriptsize, align=left]
      {slower,\\higher\\fidelity} (6.9,1.0);
    \draw[slflow] (-6.9,1.0) -- node[left, font=\scriptsize, align=right]
      {faster,\\more\\numerous} (-6.9,-4.0);
  \end{tikzpicture}%
  }
  \caption[The SmartLoad test pyramid]{The SmartLoad test pyramid, from a broad base of pure-function unit tests up to a single load-balancer-adapter conformance suite at the apex.}
  \label{fig:imc:pyramid}
\end{figure}

\begin{table}[H]
  \centering
  \caption[Per-module unit-suite test counts]{Per-module unit-suite test counts at the audited commit, the fast pure-Python suites that form the base of the pyramid.}
  \label{tab:imc:layers}
  \small
  \begin{tabular}{@{}l r l r@{}}
    \toprule
    \textbf{Module} & \textbf{Tests} & \textbf{Module} & \textbf{Tests} \\
    \midrule
    rl-engine         & 154 & policy-manager        & 65 \\
    autoscaler        & 104 & demo-ui               & 52 \\
    shared            & 103 & operator-ui           & 49 \\
    anomaly-detector  & 97  & policy-manager (more) & 38 \\
    lb-sidecar        & 79  & lb-otel-shipper       & 24 \\
    forecasting       & 53  & telemetry             & 23 \\
    lb-sidecar (conf) & 20  & resource-collector    & 20 \\
    load-balancer (conf) & 20 & experiments         & 11 \\
    \bottomrule
  \end{tabular}
\end{table}

\paragraph{Unit layer.} Each decision engine is a pure function of its inputs. The
threshold rule, the Isolation Forest scorer, the moving-average and ARIMA forecasters,
the PPO argmax-dominant weighting, and the autoscaler's dead-band arithmetic all reduce
to deterministic transforms, so each can be checked against a hand-computed
expectation. These tests run with \texttt{pytest tests/unit/} in a few seconds and gate
every commit, so a change that shifts a verdict boundary fails right away.

\paragraph{Integration layer.} These tests assume \texttt{docker compose up -d} and
check the contract between a pair of services on the wire. A \texttt{stack\_ready}
fixture in \texttt{tests/integration/conftest.py} blocks until every service reports
\texttt{/health} before any test runs. They cover telemetry OTLP ingress and database
persistence, per-request latency fidelity on live traffic (\texttt{STDDEV(latency) >
0}), the autoscaler's cooldown and reactive fallback, the dynamic-pool
\texttt{provision()}/\texttt{decommission()} lifecycle against the live Docker daemon,
and the four-channel closed-loop dispatch.

\paragraph{End-to-end layer.} These tests drive a whole feature through the
customer-facing SDK rather than an internal seam. The adaptive-bench compressed run,
for example, drives the five-phase load shape under a 60-second CI budget and checks
that all eight raw artefacts emerge with the correct schema.

\paragraph{Conformance and structural lints.} A single conformance suite
(\texttt{tests/\allowbreak conformance/\allowbreak lb\_\allowbreak adapter/}) pins the load-balancer adapter
interface. It currently exercises only the implemented \texttt{NginxAdapter} with an
in-test mock; the ALB, Envoy, and HAProxy adapters are stubs that raise
\texttt{NotImplementedError}, so the suite fixes the contract a future adapter would
have to satisfy rather than demonstrating that a second back end already meets it.
Three structural lints
(\texttt{scripts/\allowbreak lint-\allowbreak structure.py}, \texttt{lint-\allowbreak redis-\allowbreak channels.py},
\texttt{lint-\allowbreak openapi.py}) protect conventions rather than behaviour.
\Cref{tab:imc:lints} states what each enforces and where drift would otherwise hide.

\begin{table}[H]
  \centering
  \caption[The three structural lints]{The three structural lints, each turning a repository convention into a mechanical gate.}
  \label{tab:imc:lints}
  \small
  \begin{tabularx}{\linewidth}{@{}l Y Y@{}}
    \toprule
    \textbf{Lint} & \textbf{Invariant enforced} & \textbf{Drift it catches} \\
    \midrule
    \texttt{lint-structure.py}      & Every \texttt{tests/e2e/<feature>/} has a sibling \texttt{docs/features/<feature>.md} and \texttt{examples/\allowbreak scenarios/\allowbreak <feature>.py}, and every service folder carries a \texttt{README.md} & A feature shipped without its specification document, scenario, or end-to-end coverage \\
    \texttt{lint-redis-channels.py} & Every \texttt{smartload.*} channel string found in \texttt{services/} appears in \texttt{docs/redis-channels.md} & A control-bus channel published or read in code but never recorded in the registry \\
    \texttt{lint-openapi.py}        & Every Flask \texttt{/api/v1/} route in \texttt{services/} appears in \texttt{docs/\allowbreak openapi/\allowbreak smartload-v1.yaml} & A live HTTP route that the published interface contract does not document \\
    \bottomrule
  \end{tabularx}
\end{table}

\paragraph{Slow tests and the acceptance pattern.} A few live-stack tests have natural
multi-minute runtimes, such as a forecast-to-scale-out cycle or a model-calibrated
detector reaching its verdict, so they carry \texttt{@pytest.mark.slow}: the
continuous-integration \texttt{compose-test} job filters them out by default and
operators run them on demand. Beyond the layering, every product task ships \emph{two}
test artefacts, a pure-Python unit test and a live-stack acceptance test whose module
docstring cites the specification section it maps to. That citation links the
executable test to the written contract, so a reviewer can jump straight to the clause
it enforces.

\section{Evaluation Methodology}\label{imc:methodology}

The test suite answers ``does the feature work?''; the benchmarks answer ``how well,
and is it holding steady?''. SmartLoad's evaluation is built around four research
questions, three exercised by the custom harnesses inventoried in
\cref{tab:imc:harnesses}: the \emph{adaptive-bench} measures the closed
forecast~$\rightarrow$ autoscaler~$\rightarrow$~sidecar loop (RQ4); the
\emph{adaptive-advantage routing bench} compares plain NGINX round-robin against the
full decision plane on an identical, equal-capacity stack (RQ1, RQ3); and the
\emph{anomaly-engine-bench} compares the threshold rule against the trained Isolation
Forest engine across a feature grid. The adaptive-advantage bench is the canonical
routing comparison; it supersedes the earlier \emph{baseline-vs-smartload} bench, whose
50-user closed loop self-throttled to a tie that later analysis attributed to a
benchmark artefact rather than a real equivalence (\cref{imc:routingbench}). All three share one discipline: every run writes a
\texttt{MANIFEST.json} recording the timestamp, git SHA, working-tree state, and every
knob value, so a third party with the same checkout and Docker host can reproduce the
figures.

\begin{table}[H]
  \centering
  \caption[Benchmark-harness inventory]{Benchmark-harness inventory: what each harness drives, what it measures, and the research question it answers.}
  \label{tab:imc:harnesses}
  \small
  \begin{tabularx}{\linewidth}{@{}l Y l@{}}
    \toprule
    \textbf{Harness} & \textbf{What it measures} & \textbf{RQ} \\
    \midrule
    Adaptive-bench           & Closed forecast-to-scale loop under a five-phase Locust load: pool trajectory, time-to-react, upstream rewrites, anomaly recovery & RQ4 \\
    Adaptive-advantage (5v5) & Plain round-robin against the full decision plane on one equal-capacity stack across the five-phase load: per-phase and aggregate error rate and p99 tail, driven past the queue knee so a hidden bad backend sheds. Canonical; supersedes baseline-vs-smartload & RQ1, RQ3 \\
    Anomaly-engine-bench     & Threshold rule against Isolation Forest over a $16\times16$ latency-by-error-rate grid: per-cell verdict agreement and direction of divergence & --- \\
    \bottomrule
  \end{tabularx}
\end{table}

\subsection{The adaptive-bench harness}\label{imc:adaptivebench}

The adaptive-bench drives the live stack through the canonical five-phase
Locust~\cite{locust} load shape that exercises both forecast-driven scale-out and
anomaly-driven reroute in a single repeatable command. The phases, summarised in
\cref{tab:imc:phases-plan}, are: \texttt{A\_ramp} (a ramp to the steady load that sets a
forecast baseline); \texttt{B\_degrade} (a hidden backend is degraded so it sheds
\texttt{503} organically, the exclusion case round-robin cannot solve);
\texttt{C\_spike} (a sudden uniform spike the forecaster should anticipate, exercising
autoscaling); \texttt{D\_slow} (a second backend turns slow but keeps returning
\texttt{200 OK}, the latency-reroute case); and \texttt{E\_tail} (the degraded backends
recover and the system settles). An
asyncio orchestrator runs the load shape alongside three collectors: a
Prometheus~\cite{prometheus} scraper at 1\,Hz, a server-sent-events collector tailing
the decision envelopes, and an upstream watcher polling \texttt{upstream.conf} every two
seconds. A phase-D injector then drives a target backend's \texttt{/\_admin/delay}
endpoint and posts a manual isolation event.

\begin{table}[H]
  \centering
  \caption[Canonical five-phase load shape]{The canonical five-phase load shape (planned windows and events), shared by the adaptive-bench and the adaptive-advantage routing comparison.}
  \label{tab:imc:phases-plan}
  \small
  \begin{tabularx}{\linewidth}{@{}l l Y@{}}
    \toprule
    \textbf{Phase} & \textbf{Window} & \textbf{Event / what it exercises} \\
    \midrule
    \texttt{A\_ramp}    & 0--60\,s    & Ramp to steady users; sets the forecast baseline \\
    \texttt{B\_degrade} & 60--180\,s  & A hidden backend is degraded and sheds \texttt{503} organically; exclusion vs un-ejectable round-robin \\
    \texttt{C\_spike}   & 180--240\,s & Sudden uniform spike (degraded backend recovered); exercises autoscaling \\
    \texttt{D\_slow}    & 240--360\,s & A second backend turns slow but keeps returning \texttt{200 OK}; latency reroute \\
    \texttt{E\_tail}    & 360--420\,s & Degraded backends recover; the system settles \\
    \bottomrule
  \end{tabularx}
\end{table}

The join pipeline (\texttt{join\_run.py}) fuses the eight raw artefacts into a
per-second timeline (\texttt{run.parquet} plus per-channel parquets), and the plotter
(\texttt{plot\_results.py}) renders the four figures in \cref{imc:results}. To move
from single-run point estimates to publishable statistics, the harness accepts a
\texttt{-{}-runs N} flag (default five) that launches $N$ independently seeded runs; a
shared \texttt{aggregate\_runs.py} then computes, per phase and metric, the mean and a
95\% confidence interval using Student's $t$ distribution with $N-1$ degrees of
freedom, drawing confidence bands on the plots. The seed fixes only the
load-generation jitter, so the leftover run-to-run variance from cold caches,
just-in-time warm-up, and container start-up is exactly what the interval measures.

\subsection{The anomaly-engine-bench harness}\label{imc:anomalybench}

The anomaly-engine-bench is an investigation tool rather than a load test. It sweeps a
two-dimensional feature grid, with per-request latency across $[10, 600]\,\text{ms}$
and error rate across $[0, 0.2]$, asking both the threshold rule and the Isolation
Forest engine for a verdict at each cell. The output is an agreement matrix: a heatmap
of where the two engines concur and where they diverge, with the direction of any
divergence (whether the model is more cautious than the rule, or the rule catches
something the model treats as normal). The harness is deliberately decoupled from the
running cluster, calling the engines as pure functions over synthetic features so the
model's decision surface is isolated from confounds in the live telemetry path. As
\cref{imc:results} shows, it confirmed that the trained detector reproduces the
deterministic rule while adding extra coverage.

\subsection{The adaptive-advantage routing bench}\label{imc:routingbench}

The routing bench answers the comparison questions directly: it runs the same
docker-compose stack twice, once with the decision plane silent and once with it fully
engaged, and contrasts the two latency and error distributions. It supersedes an
earlier \texttt{baseline-vs-smartload} bench, whose 50-user closed loop tied for a
reason worth stating, because the same trap can silently void any load comparison. A
closed-loop generator holds a fixed population of virtual users, each issuing one
request, waiting for the response, then pausing briefly before the next, so the number
of requests in flight is capped by the user count; Locust drives exactly this model.
When a backend slows, every user routed to it blocks on the slow response and stops
generating new load, so the offered rate falls on its own and the slowdown converts into
queue-wait latency without the in-flight count ever exceeding the user population. At
50 users that population sat below the per-backend queue depth, so a degraded backend
could only add latency and could never overflow its queue; the slowdown hid in the
worst-case maximum rather than separating the two sides, and the systems tied not because
they are equivalent but because the load profile concealed the one failure mode where
they differ. That tie is a closed-loop self-throttle artefact, not a real equivalence,
and is reported as such. The
adaptive-advantage bench removes the artefact by holding the two sides at \emph{equal
capacity} and driving load past the knee. Its design principle is to minimise hidden
variables: the compose stack, backend pool image, NGINX configuration, test-backend
code, and operator-UI back-end-for-front-end are byte-for-byte identical across both
sides; only the five environment flags in \cref{tab:imc:protocol} differ, and those five
are exactly the switch that turns SmartLoad's intelligence on or off. The recorded
environment profiles live in \texttt{experiments/adaptive-advantage/env/}.

On the baseline side the services start up and expose \texttt{/health} but never enter
their inference cycle, and the lb-sidecar does not subscribe to Redis, so
\texttt{upstream.conf} stays at the static round-robin configuration. On the SmartLoad
side the same services run their full inference cycle and the sidecar rewrites
\texttt{upstream.conf} and reloads NGINX in response to the published envelopes. The
load-balancer container does not read these flags and is not restarted between sides,
so the only moving part is the decision plane.

\begin{table}[H]
  \centering
  \caption[Comparison protocol: same stack, only flags differ]{Comparison protocol, where the two environment profiles differ in exactly five variables and everything else is held identical.}
  \label{tab:imc:protocol}
  \small
  \begin{tabular}{@{}l c c@{}}
    \toprule
    \textbf{Variable} & \texttt{baseline.env} & \texttt{smartload.env} \\
    \midrule
    \texttt{ANOMALY\_RUNLOOP\_ENABLED}    & \texttt{false}  & \texttt{true}   \\
    \texttt{FORECAST\_RUNLOOP\_ENABLED}   & \texttt{false}  & \texttt{true}   \\
    \texttt{RL\_RUNLOOP\_ENABLED}         & \texttt{false}  & \texttt{true}   \\
    \texttt{LB\_SIDECAR\_RUNLOOP\_ENABLED}& \texttt{false}  & \texttt{true}   \\
    \texttt{RL\_MODE}                     & \texttt{shadow} & \texttt{active} \\
    \bottomrule
  \end{tabular}
\end{table}

The workload is the canonical five-phase Locust shape of \cref{tab:imc:phases-plan},
run at equal capacity so the comparison isolates routing and detection intelligence
rather than a capacity advantage. Three methodology choices make the result honest.
First, load is driven \emph{past the queue knee} (\texttt{STEADY\_USERS}$=70 >
\texttt{QUEUE\_MAX}=64$), so a severely degraded backend's queue overflows and sheds
\texttt{503}; with NGINX configured \texttt{max\_fails=0}, plain round-robin can never
eject it and keeps routing a fifth of its traffic onto it, while SmartLoad's error
channel detects it and the sidecar pulls it out. The \texttt{max\_fails=0} setting is a
deliberate, audited choice rather than a handicap planted to flatter SmartLoad: it
disables NGINX's passive ejection so the degraded backend keeps emitting honest
\texttt{503} backpressure instead of cascading into passive-ejection \texttt{502}s, which
isolates the variable under study, namely detection-and-exclusion intelligence against a
policy that structurally cannot exclude. Second, anomalies are injected
\emph{organically} through \texttt{/\_admin/delay} with no manual isolation hint, so the
detector must find them itself. Third, both sides are pinned to five backends
(\texttt{MIN\_BACKENDS}=\texttt{MAX\_BACKENDS}=5), removing the autoscaler-flap and
stale-\texttt{down} confounds so the 5v5 A/B measures pure routing and exclusion; between
the two sides the orchestrator also rewrites the upstream to a clean, all-up five-backend
pool before the decision plane is recreated, so a backend left \texttt{down} by the prior
side cannot carry into the next as lost capacity. This equal-capacity configuration is
the scenario the headline claim rests on, because no win can be attributed to extra
hardware; a companion full-system scenario (\texttt{MAX\_BACKENDS}=10 against a static
five) instead lets SmartLoad scale out and measures total real-world value, adaptation
plus elastic capacity, exercising the autoscaling path the pinned run deliberately holds
constant. The slow-but-healthy \texttt{D\_slow} backend that keeps returning
\texttt{200 OK} is the companion case NGINX's passive \texttt{max\_fails} check also
cannot solve. The finding this bench produces is discussed in \cref{imc:theory} and
quantified in \cref{tab:imc:phases-obs}.

\paragraph{The queue-knee arithmetic.} The threshold the workload must cross is set by
the test backend itself, which is not a constant-delay echo server but a small bounded
finite-server queue (an M/G/c model in queueing terms). Two parameters define it: two
concurrent service slots (\texttt{WORKERS}$=2$) and a bounded admission queue of depth
\texttt{QUEUE\_MAX}$=64$, so the per-backend admission ceiling is
$\texttt{WORKERS}+\texttt{QUEUE\_MAX}=66$ simultaneous requests, and the next arrival to
find the queue full is shed immediately with \texttt{503}. Under closed-loop load a
severely degraded backend holds each request far longer, so by Little's law its standing
in-flight count climbs as residence time grows; with the user population held above the
queue depth ($\texttt{STEADY\_USERS}=70>\texttt{QUEUE\_MAX}=64$) that backend accumulates
enough in-flight requests for its queue to overflow and shed organically, whereas below
the knee, in the old 50-user regime, the population is smaller than the queue and it can
never fill no matter how slow the backend gets. Driving the load just past the knee is
the single change that turns the decisive failure mode from invisible into measurable.

\subsection{Computer-based tools and statistics}\label{imc:tools}

The evaluation leans on standard, reproducible tooling. Load is driven by
Locust~\cite{locust}; system state is scraped from Prometheus~\cite{prometheus}
exposition endpoints; raw frames are stored as columnar Parquet and joined with the
\texttt{pandas}/\texttt{pyarrow} stack; plots are rendered with \texttt{matplotlib}.
Statistical reporting uses the Student's $t$ confidence interval described above for
the multi-run case, and the anomaly-engine comparison reports agreement as a simple
proportion over a fully enumerated grid, so the whole population rather than a sample
is measured and no sampling error remains to bound. The routing comparison is seeded for
reproducibility: the load stream is deterministic for a given seed and both sides of a
run share that seed, so the A/B faces an identical request stream, while successive runs
in a batch step the seed so a three-run batch exercises three independent streams. At
least three runs are reported because a single run proved able to mislead entirely; an
early batch showed a healthy first run and a collapsed second, which is why no single run
is treated as reportable.

\subsection{Threats to validity}\label{imc:threats}

Threats to validity are documented up front so a reader can weight the results
accordingly. They fall into the four standard classes, set out in
\crefrange{tab:imc:threats-internal}{tab:imc:threats-cc}. The two most material are
the load-balancer reload cost, a brief connection blip during an \texttt{upstream.conf}
rewrite that the project reports rather than filters out, and the train-versus-serve
distribution mismatch for the routing policy (\cref{imc:theory}).

\begin{table}[H]
  \centering
  \caption[Internal validity threats and controls]{Internal validity: confounds within a single run and how each is controlled.}
  \label{tab:imc:threats-internal}
  \small
  \begin{tabularx}{\linewidth}{@{}Y Y@{}}
    \toprule
    \textbf{Threat} & \textbf{Mitigation or status} \\
    \midrule
    The lb-sidecar reload causes transient connection blips that surface as max-latency spikes. & Reported honestly. The cost is part of the SmartLoad side and is included unfiltered; suppressing it would misrepresent the trade-off. \\
    NGINX's \texttt{max\_fails} passive check might trip on the baseline side if the anomaly is severe enough. & The latency-spike anomaly keeps the backend returning \texttt{200 OK}, so passive checks do not trip; the slow backend verifiably stays in rotation. \\
    The injected-degradation severity (\texttt{SEVERE\_MS}, \texttt{MODERATE\_MS}) is set by environment variables; a stale value would silently change results. & The \texttt{MANIFEST.json} captures every knob, and cross-run comparison is only treated as valid when the knobs match. \\
    Dev-stack background load (operator-UI polling, demo-ui SSE) keeps generating low-rate traffic during a run. & Both sides see the same background load, so it cancels in the delta; only the absolute numbers carry the few extra RPS. \\
    \bottomrule
  \end{tabularx}
\end{table}

\begin{table}[H]
  \centering
  \caption[External validity threats]{External validity: how far the results generalise beyond the harness.}
  \label{tab:imc:threats-external}
  \small
  \begin{tabularx}{\linewidth}{@{}Y Y@{}}
    \toprule
    \textbf{Threat} & \textbf{Mitigation or status} \\
    \midrule
    The canonical bench is equal-capacity, so it cannot reward fine-grained routing; a learned policy can only earn its keep on a heterogeneous-capacity pool, which is not yet built. & Recorded and discussed in \cref{imc:theory}. A heterogeneous-capacity testbed and retraining on a heterogeneous-latency distribution are tracked future work; until then the routing finding is scoped to equal capacity. \\
    The harness uses a single backend image, so it does not exercise non-uniform per-request CPU cost. & Acknowledged. Real-traffic shadow-mode runs are a deferred evaluation. \\
    The harness runs on a single Docker host, so cross-host network behaviour is not exercised. & Acknowledged. Kubernetes-deployment evaluations are deferred. \\
    \bottomrule
  \end{tabularx}
\end{table}

\begin{table}[H]
  \centering
  \caption[Construct and conclusion validity threats]{Construct and conclusion validity: whether the metrics measure what they claim and whether a single run supports it.}
  \label{tab:imc:threats-cc}
  \small
  \begin{tabularx}{\linewidth}{@{}l Y Y@{}}
    \toprule
    \textbf{Class} & \textbf{Threat} & \textbf{Mitigation or status} \\
    \midrule
    Construct  & ``SmartLoad beats baseline'' has no single scalar dominance criterion. & The acceptance gate reads ``clear delta on $\geq 3$ of 5 metrics''; the comparison is intentionally multi-dimensional and reported per metric. \\
    Construct  & p99 and higher tail latencies on the short two-minute profile are noisy. & Short-mode statistics are treated as validation only; the full-length profile gives tighter distributions but has not yet been re-run with a retrained model. \\
    Conclusion & A single run per side would not control for run-to-run variance from cold caches, warm-up, or container start-up. & The canonical adaptive-advantage routing bench averages three runs on a clean pinned pool, and the once-volatile \texttt{C\_spike} phase is now stable across all three; extending the same run-of-$N$ averaging with confidence intervals to every phase is a tracked next step. \\
    \bottomrule
  \end{tabularx}
\end{table}

\section{Results}\label{imc:results}

All numbers below are lifted straight from the frozen run artefacts, with nothing
interpolated.

\subsection{Adaptive-bench: the closed loop under load}\label{imc:results_adaptive}

The adaptive-bench drives the closed forecast-to-scale loop through the five-phase
load shape, recording the pool trajectory, per-phase latency and error behaviour, and
the autoscaler's scaling actions. The qualitative result is the one the design
predicts: the pool grows under the \texttt{C\_spike} surge, holds while a degraded
backend is excluded, and settles in the \texttt{E\_tail} phase without flapping, while
the controller holds error near zero in the steady phases and sheds load in a controlled
way when offered demand exceeds the host-capped pool rather than collapsing into a 502
cascade. The realised pool range and scaling-action count for this RQ4 trajectory are
being regenerated from a fresh benchmark batch and reported once it lands.
% TODO(vps-rerun): adaptive-batch RQ4 pool range + scaling-action count from fresh VPS run
\Cref{fig:imc:pool} shows the pool trajectory.

The decisive routing-and-exclusion result comes from the canonical
\emph{adaptive-advantage} bench at equal capacity (5v5, both NGINX round-robin and
SmartLoad capped at five backends, averaged over three runs on a clean pinned pool).
\Cref{tab:imc:phases-obs} reports it per phase. The headline is unambiguous: at the two
failure modes SmartLoad is built for, it decisively beats static round-robin. In
\texttt{B\_degrade} a hidden backend slow enough to shed \texttt{503} but never trip
NGINX's passive health check cuts errors from $6.53\%$ to $1.21\%$ and the p99 tail from
$56{,}000$\,ms to $223$\,ms (roughly $250\times$); in \texttt{D\_slow} a slow-but-healthy
backend cuts errors from $3.78\%$ to $0.97\%$ and the p99 tail from $14{,}000$\,ms to
$377$\,ms (roughly $37\times$). Round-robin keeps dispatching a fifth of its traffic onto
the bad backend because \texttt{max\_fails=0} leaves it un-ejectable; SmartLoad detects
it organically and excludes it. Across the whole run that exclusion advantage shows up as
roughly $5\times$ fewer errors ($1.60\%$ against $1.32\%$ overall) and an $8\times$ better
aggregate tail (p99 $4{,}867$\,ms against $623$\,ms), at equal capacity and robust across
all three runs.

Under the \texttt{C\_spike} phase the two approaches essentially tie, with round-robin
slightly ahead ($1.82\%$ errors against SmartLoad's $3.75\%$; p95 $220$\,ms against
$580$\,ms). This is the expected, honest result rather than a regression: on a
homogeneous, equal-capacity pool under uniform overload an even split is provably
optimal, so neither a learned nor a heuristic router can beat round-robin, and
SmartLoad's bounded routing skew (at most $1.33{:}1$) costs a little. SmartLoad's
measurable value is therefore anomaly-driven exclusion and capacity holding, not
fine-grained routing on a homogeneous pool.

\begin{figure}[H]
  \centering
  \bmfig{Backend pool size over the run, scaling up under the forecast burst and sustain phases and contracting again as load drops. Plot pending the host-migration benchmark batch.}
  \caption[Backend pool size over the run]{Backend pool size over the run, scaling up under the forecast burst and sustain phases and contracting again as load drops.}
  % TODO(vps-rerun): pool figure range + scaling-action count from fresh VPS run
  \label{fig:imc:pool}
\end{figure}

\begin{table}[H]
  \centering
  \caption[Adaptive-advantage per-phase results, equal capacity]{Adaptive-advantage per-phase results at equal capacity (5v5, three runs, clean pinned pool): NGINX round-robin baseline against SmartLoad, error rate and p99 tail latency.}
  \label{tab:imc:phases-obs}
  \footnotesize
  \setlength{\tabcolsep}{5pt}
  \begin{tabular}{@{}l r r r r l@{}}
    \toprule
    \textbf{Phase} & \textbf{Base err\%} & \textbf{SL err\%} & \textbf{Base p99} & \textbf{SL p99} & \textbf{Reading} \\
    \midrule
    \texttt{A\_ramp}    & 0.00\% & 0.00\% & 117\,ms     & 123\,ms & warm-up, no fault \\
    \texttt{B\_degrade} & 6.53\% & \textbf{1.21\%} & 56{,}000\,ms & \textbf{223\,ms} & exclusion win ($\sim$250$\times$ tail) \\
    \texttt{C\_spike}   & \textbf{1.82\%} & 3.75\% & 317\,ms     & 690\,ms & uniform overload; even-split RR optimal (tie) \\
    \texttt{D\_slow}    & 3.78\% & \textbf{0.97\%} & 14{,}000\,ms & \textbf{377\,ms} & latency reroute win ($\sim$37$\times$ tail) \\
    \texttt{E\_tail}    & 0.00\% & 0.00\% & 133\,ms     & 407\,ms & recovery, settle \\
    \midrule
    \textbf{Overall}    & 1.60\% & \textbf{1.32\%} & 4{,}867\,ms & \textbf{623\,ms} & $\sim$5$\times$ fewer errors, $\sim$8$\times$ better tail \\
    \bottomrule
  \end{tabular}
\end{table}

\paragraph{What makes the spike robust: an ablation.} An earlier batch of this same
bench produced a catastrophic \texttt{C\_spike} collapse, with SmartLoad's error rate
swinging as high as roughly $56\%$ between runs. That collapse was not an inherent
property of the routing plane but a coupled-failure contamination artefact: the
autoscaler flap stripped capacity mid-spike, weight concentration on the SmartLoad side
overloaded a single backend, and a benching cascade then removed it from rotation. With
those defects fixed, the spike is a small, stable trail rather than a collapse, and
\texttt{C\_spike} is now robust across all three runs (SmartLoad $3.4\%$ / $3.4\%$ /
$4.5\%$). A leave-one-out ablation (\cref{tab:imc:ablation}) decomposes which fix bought
how much of that recovery, measured as the cost on \texttt{C\_spike} error rate of
removing each fix one at a time. The anti-concentration routing clamp does the heavy
lifting: removing it alone pushes \texttt{C\_spike} from $3.3\%$ to $18.3\%$ ($+15.0$
percentage points). The equal-capacity pin is a clear second ($+1.8$), the per-side
reset a minor contributor ($+0.7$), and the absolute-overload guard is insurance that
did not need to fire in this clean run ($-0.4$). The clamp is the single most defensible
mechanism behind the recovered spike.

\begin{table}[H]
  \centering
  \caption[Spike-fix ablation: per-fix contribution]{Leave-one-out ablation of the spike-fix stack (adaptive-advantage 5v5, two runs). $\Delta$ is the cost on \texttt{C\_spike} error rate of removing that fix; positive means worse without it.}
  \label{tab:imc:ablation}
  \small
  \begin{tabular}{@{}l r r l@{}}
    \toprule
    \textbf{Fix removed} & \textbf{C\_spike err\%} & \textbf{$\Delta$ vs full} & \textbf{Reading} \\
    \midrule
    \emph{(full stack)}            & 3.3\%  & ---    & baseline round-robin here $=1.7\%$ \\
    Anti-concentration clamp       & 18.3\% & $+15.0$ & dominant; recovers the spike almost alone \\
    Equal-capacity pin             & 5.1\%  & $+1.8$  & secondary; the autoscaler flap costs capacity \\
    Per-side routing reset         & 4.0\%  & $+0.7$  & minor; stale-\texttt{down} carryover \\
    Absolute-overload guard        & 2.9\%  & $-0.4$  & insurance; did not fire in this clean run \\
    \bottomrule
  \end{tabular}
\end{table}

\Cref{fig:imc:react} plots the time-to-react: the delay between a forecast being
published and the first autoscaler action that responds to it. The forecaster emits an
envelope on a fixed cadence while the autoscaler actuates in discrete cooldown-gated
steps, so the per-forecast delay ranges from sub-second, when a forecast lands just
before an actuation, up to the cooldown-bounded maximum for the earliest ramp-phase
forecasts that must wait for the spike to cross a decision band. The measured delay
distribution for this run is regenerated from the fresh benchmark batch.
% TODO(vps-rerun): time-to-react delay distribution (min/typical/max) from fresh VPS run

\begin{figure}[H]
  \centering
  \bmfig{Time-to-react per forecast: the delay from a forecast publication to the first autoscaler action responding to it. Plot pending the host-migration benchmark batch.}
  \caption[Time-to-react per forecast]{Time-to-react per forecast: the delay from a forecast publication to the first autoscaler action responding to it.}
  \label{fig:imc:react}
\end{figure}

\Cref{fig:imc:upstream} shows the upstream timeline, the watcher's view of
\texttt{upstream.conf} over the run, and \cref{fig:imc:recovery} the anomaly-recovery
window, in which a target backend was given an injected delay at the start of phase~D
and recovered after 60 seconds. Both the anomaly POST and the recovery POST returned
\texttt{200}, confirming the injection and recovery path end-to-end.

\begin{figure}[H]
  \centering
  \bmfig{Upstream timeline: the upstream-watcher's view of the live NGINX backend list across the run. Plot pending the host-migration benchmark batch.}
  \caption[Upstream timeline across the run]{Upstream timeline: the upstream-watcher's view of the live NGINX backend list across the run.}
  \label{fig:imc:upstream}
\end{figure}

\begin{figure}[H]
  \centering
  \bmfig{Phase-D anomaly window, where a target backend receives an injected delay and is recovered after a 60-second window. Plot pending the host-migration benchmark batch.}
  \caption[Phase-D anomaly window]{Phase-D anomaly window, where a target backend receives an injected delay and is recovered after a 60-second window.}
  \label{fig:imc:recovery}
\end{figure}

\subsection{Anomaly engine: the trained detector tracks the rule}\label{imc:results_anomaly}

The anomaly-engine-bench compares the trained Isolation Forest engine against the
deterministic threshold rule across a fully enumerated feature grid. The sweep is a
calibration check, not the deployed-engine selection: the production default is the
stateful \texttt{trend\_rule} (\cref{imb:threshold}), and the threshold rule here is a
reference rather than a shipped default. The Isolation Forest model was
trained on the Server Machine Dataset~\cite{su2019omnianomaly} and scores
F1\,$=$\,\BMifFone{} on the held-out split, clearing the project's $>0.80$ acceptance
target. On the agreement sweep it covers a $16 \times 16 = 256$-cell grid and measures
a \BManAgree{} agreement rate (\BManAgreeCells{} of \BManGridCells{}) after the model
was re-calibrated in production-shape feature space. The per-engine verdict
distribution and the direction of the disagreements appear in \cref{tab:imc:anomaly}.
The remaining disagreements are \emph{all} in the value-add direction: cells where the
threshold rule says \texttt{healthy} but the model says \texttt{unhealthy}, catching
latency-variance and error-rate compounding that the static cut-off misses, with no
under-reactions, that is, no cells where the model relaxes a verdict the rule marks
unhealthy. \Cref{fig:imc:heatpost} renders the agreement matrix.

\begin{table}[H]
  \centering
  \caption[Anomaly-engine comparison over the feature grid]{Anomaly-engine comparison: threshold rule against Isolation Forest over a $16\times16=256$-cell latency-by-error-rate grid.}
  \label{tab:imc:anomaly}
  \small
  \begin{tabular}{@{}l r r@{}}
    \toprule
    \textbf{Verdict} & \textbf{Threshold} & \textbf{Isolation Forest} \\
    \midrule
    \texttt{healthy}   & \bmtodo{an-thr-healthy}   & \bmtodo{an-if-healthy}   \\
    \texttt{unhealthy} & \bmtodo{an-thr-unhealthy} & \bmtodo{an-if-unhealthy} \\
    \midrule
    \multicolumn{3}{@{}l}{\emph{Agreement: \BManAgree{} (\BManAgreeCells{}/\BManGridCells{}).}} \\
    \multicolumn{3}{@{}l}{\emph{Disagreements: all
      \texttt{healthy}$\rightarrow$\texttt{unhealthy} (value-add).}} \\
    \multicolumn{3}{@{}l}{\emph{Relaxations
      (\texttt{unhealthy}$\rightarrow$\texttt{healthy}): none.}} \\
    \bottomrule
  \end{tabular}
\end{table}

\begin{figure}[H]
  \centering
  \bmfig{Anomaly agreement heatmap, where the Isolation Forest engine agrees with the threshold rule on a high fraction of cells with no relaxed verdicts. Plot pending the host-migration benchmark batch.}
  \caption[Anomaly agreement heatmap]{Anomaly agreement heatmap, where the Isolation Forest engine agrees with the threshold rule on \BManAgree{} of cells with no relaxed verdicts.}
  \label{fig:imc:heatpost}
\end{figure}

\section{Comparison to Theoretical Predictions}\label{imc:theory}

This section sets the measured behaviours against the analytic expectations from the
design chapters. In each case the system behaved as predicted, including where the
prediction is that a learned policy ties its classical baseline.

\paragraph{Pool size versus the capacity model.} The autoscaler's design predicts a
target instance count of $n = \lceil \hat{r} / c \rceil$, where $\hat{r}$ is the
predicted request rate and $c$ the per-backend capacity, with a dead band
$[(n{-}1)c,\, nc]$ damping oscillation around the boundary and a hard cap at the
policy's \texttt{max\_backends}. The observed trajectory matches this shape: the pool
grows as the forecast burst pushes $\hat{r}$ across successive band ceilings, saturates
at the cap, and contracts symmetrically as load falls (\cref{fig:imc:pool}), with
scaling actions landing on band crossings rather than flapping. The match is structural
rather than coincidental, since the autoscaler is a direct implementation of the ceiling
rule. The realised pool range and scaling-action count for this run are regenerated from
the fresh benchmark batch.
% TODO(vps-rerun): pool range + scaling-action count tying capacity model to observed trajectory, from fresh VPS run

\paragraph{Forecast-driven scale-out and scale-in.} The design predicts that the
forecaster anticipates rising load and the autoscaler pre-warms capacity before
saturation, then sheds it as load falls. Both halves occurred as predicted: scale-out
arrived during the \texttt{C\_spike} surge, and scale-in followed cleanly as load fell
in the \texttt{E\_tail} phase, confirming the loop closes end-to-end.

\paragraph{Anomaly agreement versus the deterministic rule.} The Isolation Forest
engine was expected to reproduce the threshold rule's verdicts where the rule is
clearly right, while adding sensitivity where a single static cut-off is too blunt.
The \BManAgree{} agreement with zero relaxed verdicts is that prediction realised: the
model is a strict superset of the threshold rule's caution, and the residual
disagreements all lie in the direction of catching additional signal.

\paragraph{Learned routing versus round-robin.} On the \emph{open-loop} evaluation
the trained PPO routing policy \emph{matches} NGINX round-robin: both record an
identical \texttt{mean\_reward} of \BMrlPpoReward{} and a zero SLO-violation rate,
because the policy was trained on traces whose per-backend latencies are homogeneous
at around 9.7\,ms, all well below the 200\,ms SLO. On a homogeneous pool the optimal
routing policy \emph{is} uniform spread, since there is no faster backend to prefer, so
the policy correctly learned round-robin equivalence. The closed-loop
adaptive-advantage bench confirms the same conclusion from the other direction. On its
\texttt{C\_spike} phase, a uniform overload of an equal-capacity pool, an even split is
provably optimal, so neither the deployed \texttt{monotone} heuristic nor a learned
policy can beat round-robin: round-robin holds a marginally lower error rate ($1.82\%$
against SmartLoad's $3.75\%$), and SmartLoad's bounded routing skew costs a little. The
measurable advantage SmartLoad does carry on that bench comes not from routing but from
anomaly-driven \emph{exclusion} of the un-ejectable bad backends in \texttt{B\_degrade}
and \texttt{D\_slow} (\cref{tab:imc:phases-obs}). The deployed router is therefore the
deterministic \texttt{monotone} heuristic, kept because it ties round-robin and passes
the monotonicity probe a sane router must, while the trained PPO policy is pinned to
\texttt{shadow} mode as a dormant comparator that produces a parallel, auditable
decision stream while earning nothing on the live path. The binding constraint is named
plainly: a learned router could only earn its keep on a \emph{heterogeneous}-capacity
pool, where backends genuinely differ, which the canonical equal-capacity bench does not
exercise; building that testbed and retraining on a heterogeneous-latency distribution
under a closed-loop simulator is future work before a learned policy could be promoted
to \texttt{active}.

\paragraph{The cost of an active rewrite.} The routing bench also surfaces a cost the
project reports rather than hides. On the steady phases of the equal-capacity run the
per-phase p95 figures track closely between the two sides, but the single maximum
latency across the run is higher on the SmartLoad side. That gap is not a routing
regression; it is the lb-sidecar's NGINX-reload cost during the exclusion window, the
brief connection blip an \texttt{upstream.conf} rewrite incurs. It is included unfiltered because suppressing it
would misrepresent the trade-off between a static and a reactive configuration. The
exact maximum-latency figures for each side are regenerated from the fresh benchmark
batch.
% TODO(vps-rerun): baseline vs SmartLoad max-latency figures (reload-cost spike) from fresh VPS run

\paragraph{A broader head-to-head.} A planned, broader benchmark would compare
SmartLoad against plain round-robin across additional, more heterogeneous workloads,
quantifying end-to-end latency and failure-rate behaviour at full run length with 95\%
confidence intervals. This is a deliberate next step, building on the harness and
statistics already in place.

\section{Implementation Timeline and Work Breakdown Structure}\label{imc:wbs}

The project ran across the 2025--2026 academic year over six development phases.
\Cref{tab:imc:wbs} summarises each task with its status. The first four phases, from
requirements and architecture through the decision-plane engines to the control plane
and operator console, are complete. The benchmarking and evaluation phase is largely
complete, with the harnesses shipped and producing committed results and the broader
head-to-head set as a planned next step. The hardening and writing phase is in
progress.

\begin{table}[H]
  \centering
  \caption[Work breakdown structure phases and status]{Work breakdown structure: phases, principal deliverables, and status.}
  \label{tab:imc:wbs}
  \small
  \begin{tabularx}{\linewidth}{@{}l Y l@{}}
    \toprule
    \textbf{Phase} & \textbf{Principal deliverables} & \textbf{Status} \\
    \midrule
    Requirements \& architecture & Problem analysis, research questions, service decomposition & Complete \\
    Data plane \& telemetry & Load balancer, OTel shipper, telemetry ingest, TimescaleDB spine & Complete \\
    Decision-plane engines & Threshold + Isolation Forest, forecasters, PPO router, autoscaler & Complete \\
    Control plane \& operator UI & Policy manager, lb-sidecar, four-channel dispatch, Dev Console & Complete \\
    Benchmarks \& evaluation & Adaptive-bench, anomaly-engine-bench, multi-run CIs & In progress (broader head-to-head planned) \\
    Hardening \& writing & Calibration, structural lints, report & In progress \\
    \bottomrule
  \end{tabularx}
\end{table}

\section{Standards and Best Practices in the Implementation}\label{imc:standards}

SmartLoad's engineering rests on widely adopted standards rather than bespoke
conventions, to keep the system interoperable and its claims auditable.

\begin{itemize}
  \item \textbf{OpenAPI-linted HTTP contract.} Every service's REST surface is
        described in an OpenAPI~\cite{openapi} document, and \texttt{lint-openapi.py}
        fails the build when a route exists in code but not in the contract, so the
        published and running interfaces cannot quietly drift apart.
  \item \textbf{OTLP telemetry.} The load-balancer shipper emits per-request
        telemetry as OpenTelemetry~\cite{opentelemetry} OTLP/HTTP-JSON through a
        collector, the vendor-neutral observability standard, rather than a
        proprietary wire format.
  \item \textbf{Prometheus exposition.} Every service exposes a Prometheus
        \cite{prometheus} \texttt{/metrics} endpoint in the standard exposition
        format, which the benchmark collectors and operator dashboards scrape without
        service-specific adapters.
  \item \textbf{The engine-authoring contract.} New decision engines implement a
        single documented interface (the pluggable-engine pattern of
        \cref{ch:impl_decision}), so a threshold rule, an Isolation Forest, or a future
        model are interchangeable behind the same contract.
  \item \textbf{The per-task acceptance-test pattern.} Each task ships a unit test and
        a live-stack acceptance test whose docstring cites the specification clause it
        enforces, making the spec-to-test link explicit.
  \item \textbf{Structural CI lints.} Three lints protect the repository's own
        invariants, namely feature-to-test coverage, Redis channel publish/subscribe
        symmetry, and OpenAPI completeness, catching documentation and contract drift
        mechanically rather than by reviewer vigilance.
  \item \textbf{Semantic versioning.} Releases follow a disciplined cadence with a
        per-work-unit suffix bump, and each release row in the project's source of
        truth records its commit SHA, so any reported number traces to the exact code
        that produced it.
\end{itemize}

Taken together, these practices make SmartLoad's results not just measured but
\emph{reproducible}: a reader with the recorded git SHA, the same Docker host, and the
verbatim environment profiles can re-run a harness and get the same figures.
